% !TeX program = xelatex
\documentclass[11pt,a4paper]{article}
\usepackage{fontspec}
\setmainfont{Times New Roman}
\usepackage[russian]{babel}
\usepackage{amsmath,amssymb,tikz,enumitem}
\usepackage[left=14mm,right=14mm,top=11mm,bottom=12mm]{geometry}
\pagestyle{empty}
\setlength{\parindent}{0pt}
\setlength{\parskip}{0pt}
\linespread{1.03}
\newcommand{\heading}[1]{%
  \begin{center}
    {\large\bfseries Преобразования. Ответы}\\[1mm]
    27 октября 2026 г. \hfill Вариант #1
  \end{center}\vspace{-2mm}}
\newcommand{\ans}[2]{\par\addvspace{2mm}\noindent\textbf{#1.} #2\par}
\newcommand{\axes}[4]{%
  \draw[step=1,gray!40,line width=0.25pt] (#1,#3) grid (#2,#4);
  \draw[black!75,line width=0.6pt] (#1,0)--({#2+0.2},0);
  \draw[black!75,line width=0.6pt] (0,#3)--(0,{#4+0.2});
  \foreach \a in {#1,...,#2}{\ifnum\a=0\else
    \node[below,inner sep=1pt,font=\scriptsize,fill=white] at (\a,0) {$\a$};\fi}
  \foreach \b in {#3,...,#4}{\ifnum\b=0\else
    \node[left,inner sep=1pt,font=\scriptsize,fill=white] at (0,\b) {$\b$};\fi}
  \node[below left,inner sep=1pt,font=\scriptsize,fill=white] at (0,0) {$0$};
  \node[right,inner sep=1pt,font=\scriptsize] at ({#2+0.25},0) {$x$};
  \node[above,inner sep=1pt,font=\scriptsize] at (0,{#4+0.25}) {$y$};}
\newcommand{\graphrule}[1]{%
  \par\smallskip{\small #1}\par}
\newcommand{\ways}[1]{%
  \begin{itemize}[nosep,leftmargin=4mm,topsep=0.5mm]#1\end{itemize}}
\begin{document}
\heading{1}
\begin{minipage}[t]{0.48\linewidth}
\textbf{1.} \(f(x)=\sqrt{x+2}-1\).
\[
  \boxed{y=-2\sqrt{x+1}+3},\qquad x\geqslant-1.
\]
Переход точки: \((t;u)\mapsto(t+1;-2u+1)\).
\par\smallskip\centering
\begin{tikzpicture}[x=0.46cm,y=0.46cm]
  \axes{-2}{9}{-4}{4}
  \begin{scope}
    \clip (-2,-4) rectangle (9,4);
    \draw[thick,smooth,domain=0:4.5,samples=201,variable=\t]
      plot ({\t*\t-1},{3-2*\t});
  \end{scope}
\end{tikzpicture}
\graphrule{Точки для построения: \((-1;3)\), \((0;1)\), \((3;-1)\), \((8;-3)\).}
\end{minipage}\hfill
\begin{minipage}[t]{0.48\linewidth}
\textbf{2.} \(y=f(|x+4|)\):
\[
y=\begin{cases}
1-|x+4|,& |x+4|\leqslant1,\\
(|x+4|-1)^2,& |x+4|\geqslant1.
\end{cases}
\]
Ось симметрии: \(x=-4\).
\par\smallskip\centering
\begin{tikzpicture}[x=0.53cm,y=0.53cm]
  \axes{-7}{1}{-1}{5}
  \begin{scope}
    \clip (-7,-1) rectangle (1,5);
    \draw[thick,smooth,domain=-8:-5,samples=121]
      plot (\x,{(\x+5)^2});
    \draw[thick] (-5,0)--(-4,1)--(-3,0);
    \draw[thick,smooth,domain=-3:1,samples=161]
      plot (\x,{(\x+3)^2});
  \end{scope}
\end{tikzpicture}
\graphrule{Точки: \((-6;1)\), \((-5;0)\), \((-4;1)\), \((-3;0)\), \((-2;1)\).}
\end{minipage}

\par\addvspace{2mm}{В №\,2 сохраняем часть графика \(f\) при \(x\geqslant0\), отражаем её относительно \(Oy\), затем сдвигаем весь полученный график влево на \(4\).}

\ans{3}{\(\displaystyle \boxed{y=\frac12\sqrt{x+3}}\).}

\ans{4}{\(\boxed{g(x)=-f(x-2)+2}\).
Сдвиг вправо на \(2\), отражение относительно \(Ox\), сдвиг вверх на \(2\).
Равносильный ответ: \(g(x)=-f(-x)+2\).}

\ans{5}{\(\boxed{g(x)=-2f(x-2)+1}\).}

\ans{6}{\(\boxed{a<0,\ b=0}\).
При \(f(x)=x\) совпадение нулей требует \(b=0\).
При \(f(x)=1\) противоположность знаков требует \(a<0\).}

\ans{7}{\(\boxed{f(x)=(x+1)^2-3}\).
Проверка: \(-f(x-2)=-(x-1)^2+3\).}

\ans{8}{\(\boxed{D(g)=[-3;3],\quad E(g)=[-3;5]}\).}

\ans{9}{Например, подойдут следующие способы:
\ways{
\item растяжение от \(Ox\) в \(4\) раза, затем сдвиг вправо на \(2\);
\item сжатие к \(Oy\) в \(2\) раза, затем сдвиг вправо на \(2\):
\(y=x^2\ \longrightarrow\ y=4x^2\ \longrightarrow\ y=4(x-2)^2\);
\item сдвиг вправо на \(4\), затем сжатие к \(Oy\) в \(2\) раза:
\(y=x^2\ \longrightarrow\ y=(x-4)^2\ \longrightarrow\
y=(2x-4)^2=4(x-2)^2\).
}}

\ans{10}{Два примера перехода \((1;2)\mapsto(3;4)\):
\ways{
\item сдвиг вправо на \(2\) и вверх на \(2\): \(g(x)=f(x-2)+2\);
\item растяжение от \(Ox\) в \(2\) раза и сдвиг вправо на \(2\):
\(g(x)=2f(x-2)\).
}}
\vfill
{\footnotesize В заданиях 9 и 10 возможны и другие верные способы.}
\newpage
\heading{2}
\begin{minipage}[t]{0.48\linewidth}
\textbf{1.} \(f(x)=\sqrt{x+1}-2\).
\[
  \boxed{y=\sqrt{x+3}-5},\qquad x\geqslant-3.
\]
Переход точки: \((t;u)\mapsto(t-2;u-3)\).
\par\smallskip\centering
\begin{tikzpicture}[x=0.44cm,y=0.44cm]
  \axes{-4}{8}{-6}{1}
  \begin{scope}
    \clip (-4,-6) rectangle (8,1);
    \draw[thick,smooth,domain=0:4.5,samples=201,variable=\t]
      plot ({\t*\t-3},{\t-5});
  \end{scope}
\end{tikzpicture}
\graphrule{Точки для построения: \((-3;-5)\), \((-2;-4)\), \((1;-3)\), \((6;-2)\).}
\end{minipage}\hfill
\begin{minipage}[t]{0.48\linewidth}
\textbf{2.} \(y=f(|x-3|)\):
\[
y=\begin{cases}
1-|x-3|,& |x-3|\leqslant2,\\
(|x-3|-2)^2-1,& |x-3|\geqslant2.
\end{cases}
\]
Ось симметрии: \(x=3\).
\par\smallskip\centering
\begin{tikzpicture}[x=0.53cm,y=0.53cm]
  \axes{-1}{7}{-2}{5}
  \begin{scope}
    \clip (-1,-2) rectangle (7,5);
    \draw[thick,smooth,domain=-2:1,samples=121]
      plot (\x,{(\x-1)^2-1});
    \draw[thick] (1,-1)--(3,1)--(5,-1);
    \draw[thick,smooth,domain=5:8,samples=121]
      plot (\x,{(\x-5)^2-1});
  \end{scope}
\end{tikzpicture}
\graphrule{Точки: \((0;0)\), \((1;-1)\), \((2;0)\), \((3;1)\), \((4;0)\), \((5;-1)\), \((6;0)\).}
\end{minipage}

\par\addvspace{2mm}{В №\,2 сохраняем часть графика \(f\) при \(x\geqslant0\), отражаем её относительно \(Oy\), затем сдвигаем весь полученный график вправо на \(3\).}

\ans{3}{\(\boxed{y=3\sqrt{x-2}}\).}

\ans{4}{\(\boxed{g(x)=-f(x-2)+2}\).
Сдвиг вправо на \(2\), отражение относительно \(Ox\), сдвиг вверх на \(2\).}

\ans{5}{\(\boxed{g(x)=2f(x+3)}\).}

\ans{6}{\(\boxed{a=2,\ b=1}\).
По концам промежутков: \(-a+b=-1,\quad 2a+b=5\).}

\ans{7}{\(\boxed{f(x)=(x-1)^2-2}\).
Проверка: \(3f(x+2)=3(x+1)^2-6\).}

\ans{8}{\(\boxed{D(g)=[-2;1],\quad E(g)=[-3;2]}\).}

\ans{9}{Два способа:
\ways{
\item растяжение от \(Ox\) в \(9\) раз, затем сдвиг влево на \(1\);
\item сжатие к \(Oy\) в \(3\) раза, затем сдвиг влево на \(1\):
\(y=x^2\ \longrightarrow\ y=9x^2\ \longrightarrow\
y=9(x+1)^2\).
}}

\ans{10}{Две комбинации ровно из двух преобразований:
\ways{
\item отражение относительно \(Oy\), затем сдвиг вверх на \(2\):
\[
(-2;1)\mapsto(2;1)\mapsto(2;3),\qquad g(x)=f(-x)+2;
\]
\item сдвиг вправо на \(4\), затем растяжение от \(Ox\) в \(3\) раза:
\[
(-2;1)\mapsto(2;1)\mapsto(2;3),\qquad g(x)=3f(x-4).
\]
}}
\vfill
{\footnotesize В заданиях 9 и 10 возможны и другие верные способы.}
\end{document}
